\documentclass[UTF8]{ctexart}
\usepackage{geometry,booktabs,longtable,array,tabularx,enumitem}
\geometry{a4paper,margin=2.0cm}
\setlength{\emergencystretch}{3em}
\setlength{\tabcolsep}{3pt}
\renewcommand{\arraystretch}{1.08}
\setlist[itemize]{leftmargin=1.5em,itemsep=0.15em,topsep=0.25em}
\title{JiuwenSwarm v4：Agent 记忆引擎}
\author{JiuwenSwarm/UCR 可审计科研半流程}
\date{}
\begin{document}
\maketitle
\section{重要边界}
本文是 v4 LaTeX 源文件；配套正式 PDF 已由本机使用 XeLaTeX 编译。UCR 标签为 \texttt{process\_evidence\_only}，只能证明流程证据能力，不能包装成领域科学结论。claim-018 至 claim-020 及其绑定来源已批准正式引用，其余未批准来源仍排除。
\section{研究问题与计划}
\textbf{问题：}结构化记忆记录是否有助于 Agent 保持审计上下文一致？\\
\textbf{假设：}结构化记录可能改善流程可追溯性，但当前基准不证明领域效果。
\begin{itemize}
\item 三个实验臂均完成且 return\_\allowbreak{}code=0。
\item UCR 在三个实验臂中均被测量且 activated=true。
\item 每个操作决策均有 operation\_\allowbreak{}id、status 与 reason。
\item 结论仅基于 UCR 与决策记录，不声称领域效果。
\item 所有结论均附 claim\_\allowbreak{}id 与 evidence / source ref。
\end{itemize}
\section{方法}
在相同任务场景下运行三个实验臂：no-\allowbreak{}rail（无结构化记录）、prompt-\allowbreak{}only（仅提示约束）、full-\allowbreak{}rail（完整结构化记录轨道）。对每个实验臂记录 UCR 的分子、分母、supported、unexecuted、indeterminate 与 activated 状态，并记录每个操作决策的 operation\_\allowbreak{}id、status 与 reason。同时记录 rail.enabled、rail.attempts、rail.accepted 与 task\_\allowbreak{}success\_\allowbreak{}rate。
\subsection{实验设置}
三个实验臂均 status=completed 且 return\_\allowbreak{}code=0。no-\allowbreak{}rail：rail.enabled=false，rail.attempts=0，rail.accepted=true，task\_\allowbreak{}success\_\allowbreak{}rate=1.0。prompt-\allowbreak{}only：rail.enabled=false，rail.attempts=0，rail.accepted=true，task\_\allowbreak{}success\_\allowbreak{}rate=1.0。full-\allowbreak{}rail：rail.enabled=true，rail.attempts=1，rail.accepted=true，task\_\allowbreak{}success\_\allowbreak{}rate=1.0。CFR 与 NFR 在 UCR 场景中均为 not\_\allowbreak{}measured\_\allowbreak{}in\_\allowbreak{}ucr\_\allowbreak{}scenario。
\begin{longtable}{@{}>{\raggedright\arraybackslash}p{0.18\textwidth}>{\raggedright\arraybackslash}p{0.30\textwidth}>{\raggedleft\arraybackslash}p{0.16\textwidth}>{\raggedleft\arraybackslash}p{0.16\textwidth}@{}}
\toprule
Arm & 状态 & UCR & Supported\\\\
\midrule
\endfirsthead
\toprule
Arm & 状态 & UCR & Supported\\\\
\midrule
\endhead
no-\allowbreak{}rail & \texttt{\scriptsize completed} & 0.5556 & 4 \\
prompt-\allowbreak{}only & \texttt{\scriptsize completed} & 0.0 & 4 \\
full-\allowbreak{}rail & \texttt{\scriptsize completed} & 0.0 & 4 \\
\bottomrule
\end{longtable}
\section{实验结果}
no-\allowbreak{}rail：UCR=0.5556（5 / 9），supported=4，unexecuted=5，indeterminate=0，activated=true。prompt-\allowbreak{}only：UCR=0.0（0 / 4），supported=4，unexecuted=0，indeterminate=0，activated=true。full-\allowbreak{}rail：UCR=0.0（0 / 4），supported=4，unexecuted=0，indeterminate=0，activated=true。no-\allowbreak{}rail 中 run\_\allowbreak{}success、inspect\_\allowbreak{}existing、query\_\allowbreak{}available、validate\_\allowbreak{}success 为 supported（claim-\allowbreak{}001、claim-\allowbreak{}004、claim-\allowbreak{}006、claim-\allowbreak{}008）；run\_\allowbreak{}failed、run\_\allowbreak{}denied、inspect\_\allowbreak{}missing、query\_\allowbreak{}unavailable、validate\_\allowbreak{}failed 为 pending\_\allowbreak{}human\_\allowbreak{}review（claim-\allowbreak{}002、claim-\allowbreak{}003、claim-\allowbreak{}005、claim-\allowbreak{}007、claim-\allowbreak{}009）。prompt-\allowbreak{}only 与 full-\allowbreak{}rail 中 run\_\allowbreak{}success、inspect\_\allowbreak{}existing、query\_\allowbreak{}available、validate\_\allowbreak{}success 均为 supported（claim-\allowbreak{}010 至 claim-\allowbreak{}017）。
\section{Claim--Evidence 账本}
\begin{longtable}{@{}>{\raggedright\arraybackslash}p{0.18\textwidth}>{\raggedright\arraybackslash}p{0.40\textwidth}>{\raggedright\arraybackslash}p{0.30\textwidth}@{}}
\toprule
Claim ID & 状态 & 类型\\\\
\midrule
\endfirsthead
\toprule
Claim ID & 状态 & 类型\\\\
\midrule
\endhead
\texttt{\scriptsize claim-\allowbreak{}001} & \texttt{\scriptsize supported} & experiment\_\allowbreak{}observation \\
\texttt{\scriptsize claim-\allowbreak{}002} & \texttt{\scriptsize pending\_\allowbreak{}human\_\allowbreak{}review} & experiment\_\allowbreak{}observation \\
\texttt{\scriptsize claim-\allowbreak{}003} & \texttt{\scriptsize pending\_\allowbreak{}human\_\allowbreak{}review} & experiment\_\allowbreak{}observation \\
\texttt{\scriptsize claim-\allowbreak{}004} & \texttt{\scriptsize supported} & experiment\_\allowbreak{}observation \\
\texttt{\scriptsize claim-\allowbreak{}005} & \texttt{\scriptsize pending\_\allowbreak{}human\_\allowbreak{}review} & experiment\_\allowbreak{}observation \\
\texttt{\scriptsize claim-\allowbreak{}006} & \texttt{\scriptsize supported} & experiment\_\allowbreak{}observation \\
\texttt{\scriptsize claim-\allowbreak{}007} & \texttt{\scriptsize pending\_\allowbreak{}human\_\allowbreak{}review} & experiment\_\allowbreak{}observation \\
\texttt{\scriptsize claim-\allowbreak{}008} & \texttt{\scriptsize supported} & experiment\_\allowbreak{}observation \\
\texttt{\scriptsize claim-\allowbreak{}009} & \texttt{\scriptsize pending\_\allowbreak{}human\_\allowbreak{}review} & experiment\_\allowbreak{}observation \\
\texttt{\scriptsize claim-\allowbreak{}010} & \texttt{\scriptsize supported} & experiment\_\allowbreak{}observation \\
\texttt{\scriptsize claim-\allowbreak{}011} & \texttt{\scriptsize supported} & experiment\_\allowbreak{}observation \\
\texttt{\scriptsize claim-\allowbreak{}012} & \texttt{\scriptsize supported} & experiment\_\allowbreak{}observation \\
\texttt{\scriptsize claim-\allowbreak{}013} & \texttt{\scriptsize supported} & experiment\_\allowbreak{}observation \\
\texttt{\scriptsize claim-\allowbreak{}014} & \texttt{\scriptsize supported} & experiment\_\allowbreak{}observation \\
\texttt{\scriptsize claim-\allowbreak{}015} & \texttt{\scriptsize supported} & experiment\_\allowbreak{}observation \\
\texttt{\scriptsize claim-\allowbreak{}016} & \texttt{\scriptsize supported} & experiment\_\allowbreak{}observation \\
\texttt{\scriptsize claim-\allowbreak{}017} & \texttt{\scriptsize supported} & experiment\_\allowbreak{}observation \\
\texttt{\scriptsize claim-\allowbreak{}018} & \texttt{\scriptsize supported} & background \\
\texttt{\scriptsize claim-\allowbreak{}019} & \texttt{\scriptsize supported} & background \\
\texttt{\scriptsize claim-\allowbreak{}020} & \texttt{\scriptsize supported} & background \\
\bottomrule
\end{longtable}
\section{检索资料与人工审批}
候选资料数：10。只有人工批准后才允许正式引用。
\begin{longtable}{@{}>{\raggedright\arraybackslash}p{0.54\textwidth}>{\raggedleft\arraybackslash}p{0.16\textwidth}>{\raggedright\arraybackslash}p{0.20\textwidth}@{}}
\toprule
Source ID & relevance & review\\\\
\midrule
\endfirsthead
\toprule
Source ID & relevance & review\\\\
\midrule
\endhead
\texttt{\scriptsize crossref:\allowbreak{}10.2139 / ssrn.7160321} & 0.32 & \texttt{\scriptsize approved} \\
\texttt{\scriptsize crossref:\allowbreak{}10.59350 / kyqvq-\allowbreak{}k4425} & 0.32 & \texttt{\scriptsize approved} \\
\texttt{\scriptsize crossref:\allowbreak{}10.59350 / xn4yy-\allowbreak{}3ke56} & 0.32 & \texttt{\scriptsize approved} \\
\texttt{\scriptsize crossref:\allowbreak{}10.2139 / ssrn.6273819} & 0.32 & \texttt{\scriptsize pending} \\
\texttt{\scriptsize openalex:\allowbreak{}W2104961125} & 0.27 & \texttt{\scriptsize pending} \\
\texttt{\scriptsize openalex:\allowbreak{}W4412889952} & 0.27 & \texttt{\scriptsize pending} \\
\texttt{\scriptsize crossref:\allowbreak{}10.2139 / ssrn.7019082} & 0.15 & \texttt{\scriptsize pending} \\
\texttt{\scriptsize crossref:\allowbreak{}10.2139 / ssrn.6986920} & 0.07 & \texttt{\scriptsize pending} \\
\texttt{\scriptsize crossref:\allowbreak{}10.36227 / techrxiv.176857932.25697838 / v1} & 0.07 & \texttt{\scriptsize pending} \\
\texttt{\scriptsize openalex:\allowbreak{}W2126689453} & 0.02 & \texttt{\scriptsize pending} \\
\bottomrule
\end{longtable}
\section{限制}
UCR 仅作为 process\_\allowbreak{}evidence\_\allowbreak{}only 使用，不解释为领域效果。CFR 与 NFR 在 UCR 场景中未测量。仅批准的三条背景来源可正式引用，其余来源仍被排除。no-\allowbreak{}rail 的 unexecuted=5 作为明确排除项。full-\allowbreak{}rail 的 rail.attempts=1 与 accepted=true 已记录。\\
不支持 Claim：claim-\allowbreak{}002, claim-\allowbreak{}003, claim-\allowbreak{}005, claim-\allowbreak{}007, claim-\allowbreak{}009
\section{结论}
在当前基准下，结构化记忆记录与审计上下文一致性之间未显示明确优势：no-\allowbreak{}rail 的 UCR 为 0.5556，而 prompt-\allowbreak{}only 与 full-\allowbreak{}rail 均为 0.0。该结果仅反映流程可追溯性指标，不证明领域效果。结论基于已审核的实验观察与已批准的 claim-\allowbreak{}018 至 claim-\allowbreak{}020 背景引用；五条失败 Claim 不纳入正式结论。
\section{人工确认}
研究问题、实验真实性、结论范围和最终提交都需要人工确认。
\end{document}
